\chapter{Evaluation Framework}
\label{chap:evaluation}

This section specifies the evaluation components required for Simula methodology
compliance, as described in Sections 2.3, 3.1, and 3.3 of the research paper.

\section{Calibrated Complexity Scoring (Algorithm 3)}

Raw per-sample complexity scores are noisy due to model overconfidence.
The complexity calibrator uses batch-wise relative scoring with Elo ranking.

\subsection{Elo Rating Algorithm}

The Elo update follows the standard formula (Elo \& Sloan, 1978), adapted for
complexity comparisons. Given two examples $a$ and $b$ with ratings $R_a$ and
$R_b$, where $a$ has higher raw complexity score:

\begin{enumerate}[nosep]
    \item Compute expected outcome:
          $E_a = \frac{1}{1 + 10^{(R_b - R_a)/400}}$
    \item Update ratings with K-factor = 32:
          $R'_a = R_a + K \cdot (1 - E_a)$,
          $R'_b = R_b + K \cdot (0 - (1 - E_a))$
\end{enumerate}

\begin{lstlisting}[language=jsonx,caption={Elo computation (simula\_complexity\_calibrator.py)}]
def _update_elo(self, id_a: str, id_b: str, outcome_a: float) -> None:
    """
    Update Elo ratings for a pairwise comparison.
    
    Args:
        id_a, id_b: Example identifiers
        outcome_a: 1.0 if A wins (higher complexity), 0.0 if B wins, 0.5 for draw
        
    Standard Elo formula per Elo & Sloan (1978).
    """
    rating_a = self._scores[id_a].elo_rating
    rating_b = self._scores[id_b].elo_rating
    
    # Expected outcomes
    expected_a = 1 / (1 + math.pow(10, (rating_b - rating_a) / 400))
    expected_b = 1 - expected_a
    
    # K-factor = 32 (standard for chess, appropriate for complexity)
    K = self.config.elo_k_factor  # default 32
    
    # Update ratings
    self._scores[id_a].elo_rating = rating_a + K * (outcome_a - expected_a)
    self._scores[id_b].elo_rating = rating_b + K * ((1 - outcome_a) - expected_b)
\end{lstlisting}

\subsection{Usage}

\begin{lstlisting}[language=jsonx,caption={Complexity calibrator usage}]
from src.training.pipeline import (
    SimulaComplexityCalibrator,
    ComplexityConfig,
)

calibrator = SimulaComplexityCalibrator(
    llm_client=llm,
    config=ComplexityConfig(
        batch_size=5,   # BS=5 from Figure 10
        n_samples=5,    # N=5 from Figure 10
        enable_elo=True,
        elo_k_factor=32,  # Standard K-factor
    ),
)

# Calibrate all examples
scores = await calibrator.calibrate(examples)

# Split by complexity (Figure 7)
low, mid, high = calibrator.split_by_complexity(
    examples,
    low_percentile=0.4,
    high_percentile=0.6,
)

# Export scores
calibrator.export_scores("complexity_elo_scores.json")
\end{lstlisting}

\textbf{Key insight from paper:} High-complexity data yields +10\% accuracy on GSM8k
at 64k samples (Figure 7), but can \emph{hurt} performance when the teacher model
is weak (LEXam case).

\section{Taxonomic Coverage Evaluation}

The coverage evaluator measures Level Ratio Coverage---the proportion of unique
taxonomy nodes covered at each level:

\begin{lstlisting}[language=jsonx,caption={Coverage evaluator usage}]
from src.training.pipeline import (
    SimulaCoverageEvaluator,
    CoverageConfig,
    identify_coverage_gaps,
)

evaluator = SimulaCoverageEvaluator(
    llm_client=llm,
    config=CoverageConfig(
        max_level=4,
        sample_size=1000,
    ),
)

# Evaluate against all taxonomies
reports = await evaluator.evaluate_multiple(examples, taxonomies)

# Identify gaps
gaps = identify_coverage_gaps(reports, threshold=0.5)
print(f"Under-covered nodes: {gaps}")

# Export
evaluator.export_all_reports(reports, "coverage_report.json")
\end{lstlisting}

\section{Embedding-Based Diversity Metrics}

Following Yu et al. (2023), diversity is measured via cosine distance in embedding
space. \textbf{Global diversity} is the average pairwise distance across the dataset.
\textbf{Local diversity} is the average k-NN distance.

\begin{lstlisting}[language=jsonx,caption={Diversity analyzer usage}]
from src.training.pipeline import (
    SimulaDiversityAnalyzer,
    DiversityConfig,
)

analyzer = SimulaDiversityAnalyzer(
    config=DiversityConfig(
        embedding_model="sentence-transformers/all-MiniLM-L6-v2",
        knn_k=10,
    ),
)

report = analyzer.analyze(examples, text_field="question")
print(f"Global: {report.global_diversity:.4f}")
print(f"Local:  {report.local_diversity:.4f}")

# Identify outliers (near-duplicates or errors)
low_ids, high_ids = analyzer.identify_outliers(report)
\end{lstlisting}

\textbf{Key insight:} Global and Local diversification have \emph{additive} effects
(Figure 4). Optimizing only one is insufficient.

\section{Critic Calibration Validation}

The critic validator measures $p(y)$ and $p(y_{\text{corrupt}})$ to ensure the
double-critic is properly calibrated:

\begin{lstlisting}[language=jsonx,caption={Critic validator usage}]
from src.training.pipeline import (
    SimulaCriticValidator,
    CriticValidationConfig,
)

validator = SimulaCriticValidator(
    llm_client=llm,
    config=CriticValidationConfig(
        validation_split=0.05,
        max_samples=200,
    ),
)

report = await validator.validate(examples, schema_context)
print(f"p(y)        = {report.p_y:.3f}")
print(f"p(y_corrupt)= {report.p_y_corrupt:.3f}")
print(f"Calibrated: {report.is_well_calibrated}")
\end{lstlisting}

\textbf{Key insight:} Critic effectiveness degrades with task complexity (Figure 3b).
For effective rejection sampling, we need $p(y) > 0.8$ and $p(y_{\text{corrupt}}) < 0.3$.

\section{Student-Teacher Gap Analysis}

The paper shows that downstream performance saturates when the student bridges
approximately 83\% of the teacher-student gap (Section 4.3):

\begin{lstlisting}[language=jsonx,caption={Gap analysis}]
from src.training.pipeline import SimulaConfig

config = SimulaConfig()
config.training.teacher_accuracy = 0.70  # measured
config.training.student_baseline = 0.40  # measured

expected = config.training.expected_saturation_accuracy()
print(f"Expected saturation: {expected:.1%}")  # ~65%

# Full analysis
print(config.training.gap_analysis())
\end{lstlisting}

\section{Adaptive Features (v1.2)}

\subsection{Adaptive Complexity Ratio (Figure 7)}

The paper shows that complex data can hurt performance when the teacher model is weak
(LEXam case). The adaptive complexity feature automatically reduces the complexity
ratio when teacher accuracy is below 60\%:

\begin{lstlisting}[language=jsonx,caption={Adaptive complexity ratio}]
# In simula_config.py GenerationConfig
def get_adaptive_complexity_ratio(self, teacher_accuracy: float | None) -> float:
    """Reduce complexity for weak teachers (Figure 7)."""
    if not self.adaptive_complexity or teacher_accuracy is None:
        return self.complexity_ratio
    
    if teacher_accuracy < 0.6:
        return 0.2  # Significantly reduce
    elif teacher_accuracy < 0.7:
        return min(self.complexity_ratio, 0.35)
    return self.complexity_ratio
\end{lstlisting}

\subsection{Adaptive Critic Thresholding (Figure 3b)}

Critic effectiveness degrades with complexity. The adaptive threshold increases
the acceptance bar for high-complexity examples:

\begin{lstlisting}[language=jsonx,caption={Adaptive critic threshold}]
def get_adaptive_critic_threshold(self, complexity_score: float | None) -> float:
    """Adjust threshold based on complexity (Figure 3b)."""
    if complexity_score is None:
        return self.critic_threshold
    if complexity_score > 0.7:
        return min(0.7, self.critic_threshold + 0.2)  # More conservative
    elif complexity_score < 0.3:
        return max(0.3, self.critic_threshold - 0.2)  # More permissive
    return self.critic_threshold
\end{lstlisting}

\subsection{Automatic Coverage Gap Filling}

When coverage evaluation identifies under-represented taxonomy nodes, the system
can automatically generate additional examples to fill gaps:

\begin{lstlisting}[language=jsonx,caption={Coverage gap filling}]
from src.training.pipeline import fill_coverage_gaps, identify_coverage_gaps

# Identify gaps
gaps = identify_coverage_gaps(coverage_reports, threshold=0.5)

# Fill gaps automatically
additional = await fill_coverage_gaps(
    gaps=gaps,
    generator=generator,
    examples_per_node=10,
)
examples.extend(additional)
\end{lstlisting}

\subsection{Diversity Optimizer}

The diversity optimizer adjusts sampling weights to balance global and local diversity:

\begin{lstlisting}[language=jsonx,caption={Diversity optimizer usage}]
from src.training.pipeline import SimulaDiversityOptimizer, DiversityTargets

optimizer = SimulaDiversityOptimizer(
    targets=DiversityTargets(
        target_global=0.52,  # From Figure 4
        target_local=0.20,
    ),
)

# Analyze current diversity and get recommendations
adjustment = optimizer.analyze(diversity_report)
print(f"Global gap: {adjustment.global_diversity_gap:.3f}")
print(f"Local gap:  {adjustment.local_diversity_gap:.3f}")

# Update weights for next generation batch
new_weights = optimizer.update_weights(diversity_report, coverage_by_taxonomy)
\end{lstlisting}

\subsection{Taxonomy Quality Gate (Section B.1, Table 2)}

Generated taxonomies can be validated against expert references:

\begin{lstlisting}[language=jsonx,caption={Taxonomy quality gate}]
from src.training.pipeline import SimulaTaxonomyEvaluator, TaxonomyQualityError

evaluator = SimulaTaxonomyEvaluator(llm_client)

# Compare against expert taxonomy
report = await evaluator.compare(generated_taxonomy, "data/expert/taxonomy.json")
print(f"Completeness: {report.completeness:.1%}")  # Should be > 70%
print(f"Soundness: {report.soundness:.1%}")        # Should be > 90%

# Validate thresholds (raises TaxonomyQualityError if fails)
evaluator.validate_taxonomy(generated_taxonomy, report=report)
\end{lstlisting}

\section{Evaluation Environment Variables}

\begin{table}[H]
\centering
\caption{Evaluation Framework Environment Variables}
\label{tab:eval-env}
\begin{tabularx}{\textwidth}{lp{2cm}X}
\toprule
\textbf{Variable} & \textbf{Default} & \textbf{Description} \\
\midrule
\texttt{SIMULA\_COMPLEXITY\_BATCH\_SIZE} & 5 & BS parameter (Algorithm 3) \\
\texttt{SIMULA\_COMPLEXITY\_N\_SAMPLES} & 5 & N parameter (Algorithm 3) \\
\texttt{SIMULA\_COMPLEXITY\_ENABLE\_ELO} & true & Enable Elo calibration \\
\texttt{SIMULA\_COVERAGE\_MAX\_LEVEL} & 4 & Max taxonomy level \\
\texttt{SIMULA\_COVERAGE\_SAMPLE\_SIZE} & 1000 & Coverage sample size \\
\texttt{SIMULA\_DIVERSITY\_EMBEDDING\_MODEL} & MiniLM & Embedding model (or Gecko) \\
\texttt{SIMULA\_DIVERSITY\_KNN} & 10 & k-NN for local diversity \\
\texttt{SIMULA\_CRITIC\_VALIDATION\_SPLIT} & 0.05 & Validation fraction \\
\texttt{SIMULA\_TEACHER\_MODEL} & Qwen3.5-35B & Teacher model name \\
\texttt{SIMULA\_STUDENT\_MODEL} & Gemma-3-4B & Student model name \\
\texttt{SIMULA\_SATURATION\_RATIO} & 0.83 & Gap saturation ratio \\
\midrule
\multicolumn{3}{l}{\textbf{v1.2 Adaptive Features}} \\
\midrule
\texttt{SIMULA\_ADAPTIVE\_COMPLEXITY} & true & Auto-reduce for weak teacher \\
\texttt{SIMULA\_CRITIC\_ADAPTIVE\_THRESHOLD} & true & Complexity-aware rejection \\
\texttt{SIMULA\_CRITIC\_THRESHOLD} & 0.5 & Base critic threshold \\
\texttt{SIMULA\_TAXONOMY\_MIN\_COMPLETENESS} & 0.7 & Minimum completeness \\
\texttt{SIMULA\_TAXONOMY\_MIN\_SOUNDNESS} & 0.9 & Minimum soundness \\
\texttt{SIMULA\_DIVERSITY\_TARGET\_GLOBAL} & 0.52 & Target global diversity \\
\texttt{SIMULA\_DIVERSITY\_TARGET\_LOCAL} & 0.20 & Target local diversity \\
\texttt{SIMULA\_DIVERSITY\_TOLERANCE} & 0.05 & Diversity tolerance \\
\texttt{SIMULA\_DIVERSITY\_LEARNING\_RATE} & 0.1 & Weight adjustment rate \\
\bottomrule
\end{tabularx}
\end{table}

\section{CI Integration of Evaluation Tolerances}
\label{sec:ci-tolerances}

The reproducibility tolerances specified in the Worked Example 
(\cref{sec:worked-example}) are enforced as CI pass/fail criteria.
When running performance regression tests, the following thresholds
determine whether the test passes:

\begin{table}[H]
\centering
\caption{CI Pass/Fail Criteria for Evaluation Metrics}
\label{tab:ci-pass-fail}
\begin{tabularx}{\textwidth}{lllX}
\toprule
\textbf{Metric} & \textbf{Tolerance} & \textbf{CI Gate} & \textbf{Action on Failure} \\
\midrule
Critic rejection rate & $\pm 0.5\%$ & WARN & Log warning, continue \\
Global diversity & $\pm 0.02$ & FAIL & Block merge, investigate \\
Local diversity & $\pm 0.02$ & FAIL & Block merge, investigate \\
Coverage (Level 1) & exact & FAIL & Block merge \\
Coverage (Level 2) & $\pm 2\%$ & WARN & Log warning \\
Coverage (Level 3) & $\pm 5\%$ & WARN & Log warning \\
Mean Elo score & $\pm 50$ & INFO & Log only \\
High complexity split & $\pm 3\%$ & WARN & Log warning \\
\bottomrule
\end{tabularx}
\end{table}

\noindent
The CI script at \path{scripts/simula-regression-test.sh} implements
these checks and produces a SARIF report for integration with GitHub
Code Scanning.